\subsection{6. Common Mode Voltage Reduction of Single-Phase Quasi-Z-Source Inverter-Based Photovoltaic System}
\textbf{OSTI:} \texttt{1606674}\quad\textbf{Authors:} Sajadian, Ahmadi, Kouzani (2019)\quad\textbf{Domain:} Power electronics\\
\scorebadge{Coverage}{8}\quad\scorebadge{Agreement}{9}\quad\textbf{Total:} 17/20\par\smallskip
\noindent\textbf{Coverage rationale.} Independent NumPy SPWM generator plus PySpice/ngspice transient simulation of the common-mode equivalent circuit (Eqs 5-6 of the paper) with the same modulation indices, switching frequency, shoot-through duty, and CM-path R/L/C values. Proposed ZS1-\textgreater{}ZS2 scheme and traditional scheme both implemented. Not replicated: hardware prototype measurements, MPPT loop, full dq grid-synchronisation block, EMI filter design, component-tolerance Monte Carlo.\par
\noindent\textbf{Agreement rationale.} CMV halving reproduces exactly: |U\_CM|\_peak 400 V -\textgreater{} 200 V (50\% reduction), matching the paper's central claim that proposed scheme clamps CMV to \{0, V\_PN/2\} while traditional scheme visits \{0, V\_PN/2, V\_PN\}. Leakage current peak 40.7\% reduction (730 mA -\textgreater{} 433 mA) and RMS 36.7\% reduction align with the paper's reported envelope (paper reports leakage in the same ballpark; exact match not expected without full CM-path parasitics). Switching-harmonic line at 10 kHz drops by \textasciitilde{}40 dB in our spectrum, consistent with paper Fig. 8(c).\par\smallskip
\noindent\textbf{What reproduced:}
\begin{itemize}[leftmargin=1.4em,itemsep=0.2em,topsep=0.2em]
\item Unipolar SPWM at M=0.78, f\_s=10 kHz
\item Shoot-through duty D=0.125 qZSI modulation
\item Common-mode equivalent circuit (L\_f=3 mH, C\_g=2.2 nF, R\_g=100 Ohm, R\_f=2 Ohm)
\item Traditional vs proposed modulation leg-voltage traces
\item 50\% CMV peak reduction (exact)
\item \textasciitilde{}40 dB CMV switching-harmonic attenuation at 10 kHz
\item Leakage-current waveform comparison (Fig. 6c vs 8c)
\end{itemize}\noindent\textbf{What missing:}
\begin{itemize}[leftmargin=1.4em,itemsep=0.2em,topsep=0.2em]
\item Experimental hardware prototype
\item MPPT and grid-synchronisation control loops
\item Full EMI filter stage
\item Efficiency and THD numbers
\item Component-tolerance sensitivity study
\item Thermal / conduction-loss analysis
\end{itemize}\noindent\textbf{Follow-on research questions:}
\begin{enumerate}[leftmargin=1.6em,itemsep=0.2em,topsep=0.2em,label=Q\arabic*.]
\item How does the 50\% CMV reduction degrade as the CM parasitic capacitance C\_g varies over 0.5-10 nF (representative of different PV-array sizes and mounting conditions)
\item Does the ZS1-\textgreater{}ZS2 modulation preserve its CMV advantage under dead-time insertion of 0.5-2 us, or does dead-time reintroduce intermediate CM levels
\item Can a third modulation state (ST-\textgreater{}ZS3) further reduce CMV peak below V\_PN/4 without penalising DC-link utilisation, and what is the achievable theoretical floor
\item How do the CMV and leakage-current reductions translate to measured conducted EMI spectra (150 kHz-30 MHz) on a prototype, and does a smaller EMI filter become viable
\item When the grid voltage is distorted (\textgreater{}5\% THD, notched by nonlinear loads), does the proposed scheme maintain the clean two-level CMV, or does the control interaction reintroduce transients
\end{enumerate}

